\documentclass[12pt,a4paper]{article}

\usepackage[utf8]{inputenc}
\usepackage[T1]{fontenc}
\usepackage[spanish,es-tabla]{babel}
\usepackage{lmodern}
\usepackage{geometry}
\usepackage{longtable}
\usepackage{array}
\usepackage{booktabs}
\usepackage{hyperref}

\usepackage[backend=biber,style=ieee,sorting=none]{biblatex}
\addbibresource{referencias.bib}

\geometry{margin=2.5cm}
\hypersetup{hidelinks}
\setlength{\parindent}{0pt}
\setlength{\parskip}{6pt}

\begin{document}

\begin{center}
{\Large\textbf{Universidad Técnica Estatal de Quevedo}}\\[0.3cm]
\textbf{Facultad de Ciencias de la Computación}\\
\textbf{Carrera de Software}\\[0.4cm]
\textbf{Ingeniería de Requerimientos (ISR-401)}\\[0.6cm]
{\Large\textbf{GA-TL-02 - FORM}}\\[0.2cm]
{\large Taller de entrevista estructurada y design thinking aplicados al PFC}\\[0.6cm]
\textbf{Equipo ANY}\\
Aldas Zambrano Alan Andres\\
Naranjo Flores Anderson Jeampiere\\
Yepez Guagua Luis Antonio\\[0.4cm]
Quevedo, noviembre de 2026
\end{center}

\section{Guion de entrevista}

\section{Ensayo y consentimiento}

\section{Mapa de empatía}

\section{Definición del problema}

\section{Ajustes a las preguntas}

\section{Declaración de uso de IA}

\printbibliography[heading=bibintoc,title={Referencias bibliográficas}]

\end{document}
